\subsection{Fibered products and inverse images}

5.11.1. Let F be a Banach space, \((\mathrm{E}_{i})_{i\in I}\) a finite family of Banach spaces and \(\mathbf{f}=(f_{i})_{i\in I}\) a family of continuous linear maps \(f_{i}\) of \(E_{i}\) into F. Let E be the product of the \(E_{i}\) and let f be the continuous linear map from E into \(F^{I}\) which is the product of the \(f_{i}\). Finally, let D be the closed subspace of \(F^{I}\) formed by the \((y_{i})_{i\in I}\) such that \(y_{i}\) is independent of i (the “diagonal” of \(F^{I}\) ). The family f is said to be transversal if the continuous linear map obtained by composing f with the canonical projection of \(F^{I}\) onto the quotient \(F^{I}/D\) is surjective and if its kernel \(f^{-1}(D)\) admits a topological complement.

If the spaces \(E_{i}\) and F are all finite-dimensional, the family f is

transversal if and only if one has:

\[
\operatorname{codim} \left(\bigcap_ {i} \operatorname{Im} f _ {i}\right) = \sum_ {i} \operatorname{codim} \left(\operatorname{Im} f _ {i}\right)
\]

If \(I = \{1, 2\}\) , the pair \((f_1, f_2)\) is transversal if and only if the map

\[
f _ {1} + f _ {2}: \mathbf {E} _ {1} \oplus \mathbf {E} _ {2} \rightarrow \mathbf {F}
\]

is surjective and if its kernel admits a topological supplement (if \(E_{1}\) and \(E_{2}\) are of finite dimension, it is equivalent to say that

\[
\operatorname{Im} f _ {1} + \operatorname{Im} f _ {2} = \mathrm{F}).
\]

5.11.2. Let S be a variety, \((X_{i})_{i\in I}\) a finite family of varieties and \(\mathbf{f}=(f_{i})_{i\in I}\) a family of morphisms \(f_{i}\) of \(X_{i}\) into S. Let P be the subset of the product X of the \(X_{i}\) consisting of the points \((x_{i})_{i\in I}\) such that \(f_{i}(x_{i})\) is independent of i. Let \(x\in P\) and let \(y=f_{i}(x_{i})\in S\) . One says that the family f is transversal at the point x of P if the maps \(T_{x}(f_{i})\) form a transversal family of continuous linear maps with values in the Banach space \(T_{y}(S)\) . We say that f is transversal if it is transversal at every point of P.

If f is transversal, then P is a submanifold of X, called the fibered product of the family of \(X_{i}\) above S (relative to the family f), which we denote by \(\prod_{i\in I}X_{i}\) (or more simply \(X_{1}\times_{S}X_{2}\) when \(I=\{1,2\}\) for example). For every point \(x=(x_{i})\) of P, the tangent space \(T_{x}(\prod_{S}X_{i})\) is the subspace of \(\prod T_{x}(X_{i})\) consisting of tangent vectors \(t=(t_{i})\) such that \(T_{x_{i}}(f_{i}).t_{i}\) is independent of the index i.

If \(f_{1}: X_{1} \to Y\) is a submersion and \(f_{2}: X_{2} \to Y\) a morphism, the pair \((f_{1}, f_{2})\) is transversal.

5.11.3. ("Universal property of fibered products") Let \(\mathbf{f} = (f_i)_{i \in I}\) be a transversal family of morphisms \(f_i: X_i \to S\) and let P be the fibered product of the \(X_i\) over S; for every \(i \in I\) , one denotes by \(\pi_i\) the morphism from P into \(X_i\) obtained by restricting to P the projection of X onto \(X_i\) ; then \(f_i \circ \pi_i\) is a morphism from P to S independent of \(i\) . Let T be a variety, and let \(g_i: T \to X_i\) be morphisms of varieties such that \(f_i \circ g_i\) is a morphism from T to S independent of \(i \in I\) ; there then exists a unique morphism \(h\) from T to P such that \(g_i = \pi_i \circ h\) for all \(i \in I\) .

5.11.4. ("Associativity of the fibered product") Let \(\mathbf{f} = (f_i)_{i \in I}\) be a finite family of morphisms of varieties \(f_i: \mathbf{X}_i \to \mathbf{S}\) , and let \((\mathrm{J}_{\lambda})_{\lambda \in \Lambda}\) be a partition of I. Suppose that, for every \(\lambda\) in \(\Lambda\) , the family \(\mathbf{f}_{\lambda} = (f_i)_{i \in \mathrm{J}_{\lambda}}\) is transversal, and we denote by \(\mathrm{P}_{\lambda}\) the fiber product of this family; for every point \(x = (x_i)_{i \in \mathrm{J}_{\lambda}}\) of \(\mathrm{P}_{\lambda}\) , the element \(f_i(x_i)\) of S is independent of \(i \in \mathrm{J}_{\lambda}\) and will be denoted by \(u_{\lambda}(x)\) ;

then \(u_{\lambda}\) is a morphism from \(P_{\lambda}\) into S. For the family \(\mathbf{u} = (u_{\lambda})\) to be transversal, it is necessary and sufficient that the family \(\mathbf{f}\) be so. The canonical bijection of \(\prod_{\lambda \in \Lambda} \prod_{i \in J_{\lambda}} X_i\) onto \(\prod_{i \in I} X_i\) then gives by restriction an isomorphism from the fibered product \(\prod_{\lambda \in \Lambda} {}_{S}P_{\lambda}\) to the fibered product \(\prod_{i \in I} {}_{S}X_i\) .

5.11.5. Let S be a variety. One calls a variety over S a variety X equipped with a morphism \(\lambda: X \to S\) . Let \((X, \lambda)\) be a variety over S and \(f: S' \to S\) a morphism of varieties such that the pair \((f, \lambda)\) is transversal. We will then denote \(f^*(X)\) the variety \(S' \times_{S} X\) , endowed with the morphism \(f^*(\lambda): S' \times_{S} X\) defined by \(f^*(\lambda)(s', x) = s'\) . We say that \(f^*(X)\) is deduced from X by base change from S to \(S'\) along \(f\) . If \(\lambda\) is a submersion (resp. an immersion, a subimmersion, an étale morphism), the same holds for \(f^*(\lambda)\) .

5.11.6. Let \(f: X \to Y\) be a morphism of varieties and let W be a subvariety of Y; let \(i\) be the canonical injection of W into Y. We say that \(f\) is transverse to W at a point \(x \in f^{-1}(W)\) if the pair \((f, i)\) is transversal at the point \((x, f(x))\) of \(X \times W\) . For this, it is necessary and sufficient that the following conditions be satisfied:

\begin{enumerate}
\def\labelenumi{(\roman{enumi})}
\item
  the tangent space \(\mathbf{T}_{f(x)}(\mathbf{Y})\) is the sum of \(\mathbf{T}_{f(x)}(\mathbf{W})\) and the image of \(\mathbf{T}_{x}(f)\) ;
\item
  the inverse image \(\mathrm{T}_{x}(f)^{-1}(\mathrm{T}_{f(x)}(\mathrm{W}))\) of the tangent space to W at \(f(x)\) admits a topological complement.
\end{enumerate}

One says that \(f\) is transverse to \(W\) if it is transverse to \(W\) at every point of \(f^{-1}(W)\) .

For a morphism from X to Y to be a submersion, it suffices that it be transversal to every point of Y, and it is necessary that it be transversal to every submanifold of Y. For a finite family of morphisms \(f_{i}:X_{i}\to S\) to be transversal, it is necessary and sufficient that the morphism \(g\) from \(\prod_{i\in I}X_{i}\) to \(S^{I}\) defined by \(g((x_{i})_{i\in I})=(f_{i}(x_{i}))_{i\in I}\) be transverse to the diagonal of \(S^{I}\) .

5.11.7. Suppose that the morphism \(f:X\to Y\) is transverse to the subvariety W of Y. Then \(f^{-1}(W)\) is a submanifold of X, and for \(x\) in \(f^{-1}(W)\) the subspace of \(\mathrm{T}_{x}(\mathrm{X})\) tangent to \(f^{-1}(W)\) is the inverse image under \(\mathrm{T}_{x}(f)\) of the subspace \(\mathrm{T}_{f(x)}(\mathrm{W})\) . By passing to the quotient, the linear map \(\mathrm{T}_{x}(f)\) defines an isomorphism of topological vector spaces from the transverse space to \(f^{-1}(W)\) at \(x\) onto the transverse space to W at \(y=f(x)\) . If W has codimension \(d\) at \(y\) in Y, the submanifold \(f^{-1}(W)\) of X has codimension \(d\) at every point of \(f^{-1}(W)\) . Finally, the map \(x\mapsto(x,f(x))\) is an isomorphism of varieties from \(f^{-1}(W)\) to the fibered product \(\mathrm{X}\times_{\mathrm{Y}}\mathrm{W}\) .

5.11.8. Let \(Y_{1}\) and \(Y_{2}\) be two submanifolds of a variety X, and let \(\iota_{j}\) be the injection of \(Y_{j}\) into X. One says that \(Y_{1}\) and \(Y_{2}\) are transverse if the pair \((\iota_{1},\iota_{2})\) is transverse; this is equivalent to supposing that, for every point x of \(Y_{1}\cap Y_{2}\) , the subspaces \(T_{x}(Y_{1})\) and \(T_{x}(Y_{2})\) of \(T_{x}(X)\) have \(T_{x}(X)\) as their sum, and that their intersection admits a topological supplement in \(T_{x}(X)\) . Under these conditions, \(Y_{1}\cap Y_{2}\) is a submanifold of X and for every x in \(Y_{1}\cap Y_{2}\) , we have:

\[
\mathrm{T} _ {x} (\mathrm{Y} _ {1} \cap \mathrm{Y} _ {2}) = \mathrm{T} _ {x} (\mathrm{Y} _ {1}) \cap \mathrm{T} _ {x} (\mathrm{Y} _ {2});
\]

if moreover \(\mathbf{Y}_i\) has codimension \(d_i\) at \(x\) , then \(\mathbf{Y}_1 \cap \mathbf{Y}_2\) has codimension \(d_1 + d_2\) at \(x\) .

5.11.9. Let \(f\) and \(g\) be two morphisms from a variety X to a variety Y. If the morphism \((f,g):\mathrm{X}\to\mathrm{Y}\times\mathrm{Y}\) is transverse to the diagonal of \(\mathrm{Y}\times\mathrm{Y}\) , the subset N of X consisting of the points \(x\) such that \(f(x)=g(x)\) is a submanifold of X; it is called the kernel of the double arrow.

\[
f, g: \mathrm{X} \longrightarrow \mathrm{Y}.
\]

